\documentclass[11pt]{article}

\usepackage[a4paper,margin=2.25cm]{geometry}
\usepackage{amsmath}
\usepackage{booktabs}
\usepackage{float}
\usepackage{graphicx}
\usepackage{longtable}
\usepackage[hidelinks]{hyperref}
\usepackage{siunitx}

\IfFileExists{docs/generated/xsection_fit/fit_summary.tex}{
  \newcommand{\generateddir}{docs/generated/xsection_fit}
  \newcommand{\figuredir}{docs/figures/xsection_fit}
}{
  \newcommand{\generateddir}{generated/xsection_fit}
  \newcommand{\figuredir}{figures/xsection_fit}
}

\title{Selection, curation, and uncertainty modelling of PET-production
cross sections}
\author{J.J. G\'omez-Cadenas}
\date{\today}

\begin{document}
\maketitle

\begin{abstract}
Experimental excitation functions are used to construct cross-section
ensembles for five positron-emitter production channels relevant to
proton-therapy PET.  The input is a pinned EXFOR snapshot subjected to explicit
dataset- and point-level selection.  Positive, threshold-constrained spline
fits retain the separation between experimental campaigns and account for
reported point uncertainties and observed between-campaign dispersion.
JENDL-4.0/HE and LANL ENDF/B-VII.1 provide external comparisons and do not
enter the fit.  The nominal curve and its replicas are folded with
target-resolved proton track length from a frozen Geant4 transport and
propagated to the production and reconstructed distal activity edges.
\end{abstract}

\section{Reaction scope and role in the uncertainty calculation}

The selected residual-production channels are
\begin{equation}
\begin{split}
  &^{12}\mathrm{C}(p,x)^{11}\mathrm{C},\qquad
   ^{16}\mathrm{O}(p,x)^{15}\mathrm{O},\qquad
   ^{16}\mathrm{O}(p,x)^{11}\mathrm{C},\\
  &^{14}\mathrm{N}(p,x)^{13}\mathrm{N},\qquad
   ^{16}\mathrm{O}(p,x)^{13}\mathrm{N}.
\end{split}
\end{equation}
The fitted interval is \SIrange{5}{150}{MeV}.  This interval contains the
low-energy production that controls the distal activity fall-off and the
incident-energy range of the treatment field.  The extension beyond the
highest treatment layer constrains the upper fit boundary with measured data.
The fit uses experimental cross sections.  Geant4 subsequently supplies the
proton trajectories, energy loss, fluence and total inelastic attenuation.
Production through the five selected channels is calculated from the fitted
curves; the corresponding Geant4 residual-production cross sections and final-
state branching do not enter either the nominal result or its replicas.

\section{Experimental data and evaluated comparisons}

\subsection{EXFOR snapshot and normalized representation}

The experimental input is the EXFOR Plot File snapshot at commit
\texttt{57052199247e}, dated 29 June 2026; the complete commit identifier is
stored in the source manifest.  The selected snapshot contains 85
excitation-function series for the five channels.  Each CX4 series is converted
one-to-one into the common schema under \path{data/xsections/normalized}; the
conversion changes eV to MeV and b to mb and preserves the EXFOR accession,
reaction field, source reference, central values, and reported uncertainties.

Experimental campaigns remain separate.  The EXFOR accession prefix defines
the campaign identifier, so subentries obtained from the same irradiation
share a campaign-level contribution in the covariance model.

\subsection{Evaluated libraries}

JENDL-4.0/HE and the LANL proton evaluations distributed with ENDF/B-VII.1 are
converted to the same energy and cross-section units.  They are compared with
the resulting experimental fit.  Their values are not treated as
measurements and do not contribute to the nominal curve or its uncertainty
band.

\section{Data curation}

\subsection{Dataset-level selection}

Dataset-level curation is based on the normalized EXFOR metadata.  A series is
accepted when
\begin{enumerate}
  \item its projectile, target isotope, and residual isotope identify one of
        the five selected channels;
  \item the EXFOR reaction field identifies an absolute production cross
        section, with no relative-yield, abundance, thick-target-yield,
        monitor-ratio, or cumulative modifier;
  \item the original units are eV and b and therefore have a direct conversion
        to MeV and mb; and
  \item the normalized quantity is identified as a cross section.
\end{enumerate}
The accepted EXFOR reaction suffixes are \texttt{,,SIG},
\texttt{,,SIG,,,EXP}, and \texttt{,IND,SIG,,,EXP}.  A dataset that fails one
of these tests is assigned either \emph{excluded} or \emph{pending
verification}.  Numerical disagreement between experiments is not, by itself,
an exclusion criterion.

Six series are pending verification.  The EXFOR reaction fields for B0077.002
and B0077.003 contain the SF8 modifier \texttt{A}, which denotes a cross
section multiplied by isotopic abundance.\footnote{IAEA,
\href{https://nds.iaea.org/nrdc/nrdc_doc/iaea-nds-0208-200601.pdf}{LEXFOR,
General Quantity Modifiers and Products}.}  The stored values are consequently
not direct absolute cross sections.  These two series remain outside the fit
until the abundance factor and normalization are recovered from the original
entry and publication.

B0095.002 contains two measurements at 20.7 and \SI{21.1}{MeV}.  Its EXFOR
incident-energy field is \SI{1}{MeV}, but the present curation has not
established whether this field denotes calibration uncertainty, resolution,
or an energy interval.  The series is therefore pending and is omitted from
the nominal fit.  A named sensitivity fit includes B0095.002 together with all
nominally accepted carbon data.

Masuda E2568.002, E2568.003, and E2568.004 are relative excitation functions
normalized simultaneously to
$^{16}$O(p,x)$^{15}$O = \SI{76.8}{mb} at \SI{35}{MeV}.  Their shared external
normalization is not represented by the present covariance model.  These
series are retained in the figures at their published scale as shape
comparisons and are pending for the absolute fit.

The metadata criteria provide a reproducible initial selection.  Verification
against an original publication is recorded as a separate curation action when
the EXFOR quantity or normalization is ambiguous; the procedure does not infer
an undocumented uncertainty or normalization component.

\subsection{Point-level selection}

Within an accepted series, a point contributes to the weighted fit when
\begin{equation}
  5\leq E\leq150~\mathrm{MeV},\qquad
  E>E_{\mathrm{th}},\qquad
  \sigma>0,
\end{equation}
and EXFOR reports a positive cross-section uncertainty.  The adopted thresholds
are 17.9, 14.29, 25.17, 8.9, and \SI{5.55}{MeV}, respectively, in the channel
order given above.  For the two corrected oxygen channels, Table~\ref{tab:thresholds}
gives the lowest production route and the exact laboratory-frame threshold
calculated from AME2020 atomic masses.  The calculation uses
\begin{equation}
 E_{\mathrm{th,lab}}=
 \frac{M_f^2-(M_p+M_T)^2}{2M_T}c^2,
\end{equation}
where $M_f$ is the sum of the final-state rest masses.  The rounded values in
the fit configuration are checked against the generated calculation.

\begin{table}[H]
\centering\small
\caption{Lowest production routes and laboratory thresholds.  Atomic masses
are from the \href{https://amdc.impcas.ac.cn/web/masseval.html}{AME2020 atomic-mass evaluation}.}
\label{tab:thresholds}
\begin{tabular}{lrr}
\toprule
Lowest channel & $Q$ (MeV) & $E_{\mathrm{th,lab}}$ (MeV) \\
\midrule
$^{16}\mathrm{O}(p,d)^{15}\mathrm{O}$ & -13.439 & 14.292 \\
$^{16}\mathrm{O}(p,d+\alpha)^{11}\mathrm{C}$ & -23.658 & 25.168 \\
\bottomrule
\end{tabular}

\end{table}

Dataset and point decisions are stored in
\path{data/xsections/curation.csv} and
\path{data/xsections/point_curation.csv}.

\subsection{Curation outcome}

Table~\ref{tab:curation} summarizes the selected data.  Of the 85 EXFOR
series, 79 are accepted and six remain pending.  A total of 834 experimental
points satisfy the dataset- and point-level rules.

\begin{table}[H]
\centering\small
\caption{EXFOR curation by production channel.}
\label{tab:curation}
\begin{tabular}{lrrrr}
\toprule
Channel & Accepted & Pending & Excluded & Fit points \\
\midrule
C-12(p,x)C-11 & 27 & 3 & 0 & 94 \\
O-16(p,x)O-15 & 11 & 1 & 0 & 94 \\
O-16(p,x)C-11 & 11 & 1 & 0 & 58 \\
N-14(p,x)N-13 & 13 & 0 & 0 & 76 \\
O-16(p,x)N-13 & 17 & 1 & 0 & 512 \\
\bottomrule
\end{tabular}

\end{table}

\section{Cross-section model}

\subsection{Threshold-constrained smooth curve}

For $E>E_{\mathrm{th}}$, the cross section is represented as
\begin{equation}
  \sigma(E)=(E-E_{\mathrm{th}})^{1/2}\exp[g(E)],
  \label{eq:curve}
\end{equation}
where $g(E)$ is a cubic B-spline.  Equation~\eqref{eq:curve} imposes a zero
cross section at and below threshold and a positive cross section above it.  A
squared second-difference penalty controls the spline curvature.

The fit is performed in the transformed variable
\begin{equation}
  g_i=\log\sigma_i-\frac{1}{2}\log(E_i-E_{\mathrm{th}}).
\end{equation}
A reported fractional one-standard-deviation cross-section error $r_i$ is
mapped to the logarithmic error
\begin{equation}
  u_i=\sqrt{\log(1+r_i^2)}.
\end{equation}

\subsection{Campaign covariance and smoothing selection}

For experimental points $i$ and $j$, the covariance model is
\begin{equation}
  C_{ij}=u_i^2\delta_{ij}+s_c^2 I[c(i)=c(j)],
  \label{eq:covariance}
\end{equation}
where $c(i)$ denotes the campaign and $s_c$ is a common logarithmic
between-campaign spread.  The second term represents a coherent displacement
of a campaign rather than independent fluctuations of all its points.

The curvature penalty and $s_c$ are selected from fixed grids by
leave-one-campaign-out prediction.  Each complete campaign is predicted by a
fit to the remaining campaigns, and the selection score is the Gaussian
predictive loss including its variance term.  The fractional campaign spread
reported in the results is $\exp(s_c)-1$.

Spline knots are dense below \SI{100}{MeV}, where the threshold structure and
excitation-function maxima occur, and coarser above \SI{100}{MeV}.  The latter
region is constrained by accepted measurements through \SI{150}{MeV}; this
placement suppresses free-boundary oscillation while retaining the measured
high-energy trend.

\section{Replica construction and cross-section uncertainty}

One thousand replicas are generated with a recorded random seed.  For each
replica, every fitted point receives a draw from its reported logarithmic
cross-section error and every campaign receives one common draw with standard
deviation $s_c$.  The curvature penalty is reselected from the fixed grid, and
the varied data are refitted.

The implementation also defines within-campaign correlated draws for
separately quantified normalization errors and verified incident-energy
uncertainties.  No such numerical component is active in the current
configuration: the curated metadata contain no separate campaign
normalization value, and no CX4 energy field has yet been verified as an
uncertainty rather than an energy interval.

At each energy, the replica median defines the nominal cross section.  The
16th and 84th percentiles define the primary uncertainty interval.  Its
relative half-width is
\begin{equation}
  u_{\mathrm{rel}}(E)=
  \frac{\sigma_{84}(E)-\sigma_{16}(E)}{2\sigma_{50}(E)}.
\end{equation}

\section{Fit results}

Table~\ref{tab:summary} gives the number of campaigns and fitted points, the
fractional between-campaign spread $\exp(s_c)-1$, the median relative band
half-width in the production-bearing interval, and the cross section at the
maximum of the nominal curve.

\begin{table}[H]
\centering\small
\caption{EXFOR fit and uncertainty summary.  The peak interval is the 16th--84th
percentile interval evaluated at the maximum of the nominal curve.}
\label{tab:summary}
\setlength{\tabcolsep}{4pt}
\begin{tabular}{lrrrrrr}
\toprule
Channel & $N_{\mathrm{camp}}$ & $N_{\mathrm{pt}}$ & $e^{s_c}-1$ & Median $u$ & Peak (mb) & $\chi^2/\nu$ \\
\midrule
C-12(p,x)C-11 & 7 & 94 & 5\% & 3\% & 88.1 [86.2, 89.9] & 2.68 \\
O-16(p,x)O-15 & 7 & 94 & 11\% & 3\% & 74.9 [72.8, 77.1] & 4.57 \\
O-16(p,x)C-11 & 5 & 58 & 42\% & 5\% & 28.6 [27.3, 29.9] & 8.63 \\
N-14(p,x)N-13 & 7 & 76 & 42\% & 5\% & 29.4 [28.1, 30.9] & 4.66 \\
O-16(p,x)N-13 & 12 & 512 & 639\% & 5\% & 36.2 [34.4, 38.0] & 20.88 \\
\bottomrule
\end{tabular}

\end{table}

The reported $\chi^2/\nu$ is a lack-of-fit diagnostic evaluated with the point
errors and fitted campaign-spread covariance.  Values above unity indicate
shape-level scatter remaining after coherent campaign offsets are included.
The replica band follows the stated covariance model and is not rescaled to
force $\chi^2/\nu=1$.

For $^{12}$C(p,x)$^{11}$C, omission of B0095.002 removes the isolated
near-threshold peak from the nominal curve.  The dashed sensitivity curve in
the carbon figure adds B0095.002 to the accepted datasets.  Across the
production-bearing region its median absolute fractional difference from the
nominal fit is 0.9\%; the largest difference occurs on the steep threshold
rise.  The complete numerical comparison is stored in
\path{\generateddir/curation_sensitivity.csv}.

\begin{figure}[p]\centering
\includegraphics[width=\textwidth]{\figuredir/p_C12_x_C11.pdf}
\caption{Fit and uncertainty ensemble for
$^{12}$C(p,x)$^{11}$C.}
\end{figure}
\begin{figure}[p]\centering
\includegraphics[width=\textwidth]{\figuredir/p_O16_x_O15.pdf}
\caption{Fit and uncertainty ensemble for
$^{16}$O(p,x)$^{15}$O.}
\end{figure}
\begin{figure}[p]\centering
\includegraphics[width=\textwidth]{\figuredir/p_O16_x_C11.pdf}
\caption{Fit and uncertainty ensemble for
$^{16}$O(p,x)$^{11}$C.}
\end{figure}
\begin{figure}[p]\centering
\includegraphics[width=\textwidth]{\figuredir/p_N14_x_N13.pdf}
\caption{Fit and uncertainty ensemble for
$^{14}$N(p,x)$^{13}$N.}
\end{figure}
\begin{figure}[p]\centering
\includegraphics[width=\textwidth]{\figuredir/p_O16_x_N13.pdf}
\caption{Fit and uncertainty ensemble for
$^{16}$O(p,x)$^{13}$N.}
\end{figure}

\clearpage
\section{Comparison with evaluated libraries}

The lower-right panel of each channel figure compares the experimental fit
with JENDL-4.0/HE and LANL ENDF/B-VII.1.  The numerical comparison is evaluated
over the common interval in which the nominal fit exceeds 5\% of its maximum.
For each evaluated curve,
\path{docs/generated/xsection_fit/evaluation_comparison.csv} records the median
ratio to the experimental fit and the median and maximum absolute fractional
differences.  These comparisons diagnose evaluation dependence; they are not
samples of the experimental uncertainty distribution.

\section{Replica-distance diagnostic}

For plotting purposes, replica $r$ is assigned the distance
\begin{equation}
  D_r=\frac{1}{N}[g_r-g_{50}]^{\mathsf T}C^{-1}[g_r-g_{50}],
  \label{eq:distance}
\end{equation}
evaluated at the $N$ fitted experimental energies.  The histogram in each
figure contains all replicas, and the nine displayed curves are selected at
fixed quantiles of $D$.

Equation~\eqref{eq:distance} is a diagnostic of the replica ensemble.  Its
magnitude depends on the channel covariance, the number and placement of the
data, and the selected curvature penalty.  It is not used to rank channel
uncertainties and does not enter the propagated distal-edge uncertainty.
The curvature penalty is reselected from a discrete grid for every replica.
Consequently, replicas assigned to different penalty values can form separate
modes in the $D$ histogram.  The counts and distance range for each selected
penalty are recorded in
\path{\generateddir/distance_by_smoothing.csv}; the distance percentile is
used only to choose the nine displayed examples.

\section{Propagation to the distal activity edge}

Let $s$ denote a proton step of length $\ell_s$, representative kinetic energy
$E_s$, and statistical weight $a_s$.  If the step traverses a material with
number density $n_{t,s}$ of target isotope $t$, its expected contribution to
production channel $c$ in cross-section replica $r$ is
\begin{equation}
  \Delta N_{s,c}^{(r)}=
  a_s n_{t(c),s}\ell_s\sigma_c^{(r)}(E_s).
  \label{eq:fold-step}
\end{equation}
Summing Eq.~\eqref{eq:fold-step} over the steps in spatial bin $m$ and over the
channels producing isotope $j$ gives
\begin{equation}
  P_{j,m}^{(r)}=
  \sum_{s\in m}\sum_{c\rightarrow j}
  a_s n_{t(c),s}\ell_s\sigma_c^{(r)}(E_s).
  \label{eq:fold-profile}
\end{equation}
The nominal prediction uses the replica median $\sigma_{50}(E)$; each
uncertainty realization substitutes one complete replica curve.  All
realizations use the same stored proton steps, so differences among them arise
from the experimental cross-section ensemble.

Geant4 continues to determine proton and secondary-proton transport, energy
loss, total inelastic attenuation and dose.  Residual nuclei generated by the
Geant4 final-state model through the five selected channels are excluded from
the folded source tally, preventing double counting.  Their production is
replaced by Eqs.~\eqref{eq:fold-step} and \eqref{eq:fold-profile}.  This method
also remains defined where a Geant4 residual-production model has zero support,
because it contains no channel-cross-section denominator.

The fixed-transport approximation holds the proton fluence constant while the
selected partial production cross sections vary.  These channels contribute
only a fraction of the total proton inelastic cross section, so the associated
change in fluence attenuation is a higher-order correction.  The implementation
will record the ratio of their summed cross section to the Geant4 total
inelastic cross section as a quantitative check of this approximation.

Equation~\eqref{eq:fold-profile} directly provides isotope budgets,
production-depth profiles and the production-edge displacement.  A common
proposal sample drawn from the scored proton steps supplies production sites
for positron transport.  Nominal and replica sources use the same proposal
sites and random numbers, with probabilities proportional to
$\Delta N_{s,c}^{(r)}$, thereby preserving correlations and reducing sampling
noise.  The resulting annihilation coordinates pass through the common
beam-off acquisition, detector simulation, reconstruction and edge estimator.

The final cross-section uncertainty is obtained from the distribution of
distal-edge displacements
\begin{equation}
  b_{\mathrm{reco}}^{(r)}=
  R_{50,\mathrm{reco}}^{(r)}-R_{50,\mathrm{reco}}^{(\mathrm{nominal})},
\end{equation}
after applying the common acquisition, detector simulation, reconstruction,
and edge estimator.  The production-profile displacement is recorded
separately.  The median and 16th--84th percentiles of the paired displacement
distribution define the propagated bias and uncertainty interval.  Thus all
replicas enter the physical propagation; the value of $D_r$ has no role in
their weighting.

\section{Conclusions}

The curation retains 79 absolute EXFOR series for the five selected production
channels and assigns six series to pending verification.  The fitted ensembles
preserve reported point errors, coherent campaign variation, and uncertainty
in the smooth representation.  Folding each complete curve with a common
Geant4 proton track-length field transfers this uncertainty to isotope budgets
and spatial production profiles without using Geant4 residual-production
cross sections.  The
between-campaign dispersion is channel dependent and is largest for
$^{16}$O(p,x)$^{13}$N in the present selection.  The physical importance of
each variation will be determined by propagation to the reconstructed distal
edge rather than by its distance from the nominal cross-section curve.

\appendix

\section{Dataset-level curation audit}

Table~\ref{tab:curation-detail} lists every selected EXFOR series, its energy
coverage, dataset decision, total number of points, and number satisfying the
numerical point rules.  The latter count does not supersede the dataset
decision: points from a pending series remain outside the fit.  The full audit
table additionally records the reaction field, quantity, reference, DOI,
decision basis, and notes.

\begingroup
\small
\setlength{\LTcapwidth}{\textwidth}
\begin{longtable}{llllrr}
\toprule
\caption{Dataset-level EXFOR curation audit.}\label{tab:curation-detail}\\
Accession & Channel & Energy (MeV) & Decision & Points & Point-rule pass \\
\midrule
\endfirsthead
\toprule
Accession & Channel & Energy (MeV) & Decision & Points & Point-rule pass \\
\midrule
\endhead
A0660.018 & C-12(p,x)C-11 & 1870.00--2690.00 & accepted & 2 & 0 \\
B0076.005 & C-12(p,x)C-11 & 8.50--104.00 & accepted & 42 & 0 \\
B0077.002 & C-12(p,x)C-11 & 15.90--32.00 & pending verification & 42 & 1 \\
B0077.003 & C-12(p,x)C-11 & 93.00--340.00 & pending verification & 8 & 2 \\
B0094.003 & C-12(p,x)C-11 & 591.00--591.00 & accepted & 1 & 0 \\
B0095.002 & C-12(p,x)C-11 & 20.70--50.50 & pending verification & 9 & 9 \\
B0101.002 & C-12(p,x)C-11 & 170.00--350.00 & accepted & 10 & 0 \\
C0058.002 & C-12(p,x)C-11 & 50.70--158.60 & accepted & 17 & 16 \\
C0062.011 & C-12(p,x)C-11 & 5700.00--5700.00 & accepted & 1 & 0 \\
C0206.002 & C-12(p,x)C-11 & 288.00--383.00 & accepted & 5 & 0 \\
C0216.002 & C-12(p,x)C-11 & 2000.00--3000.00 & accepted & 6 & 0 \\
C0220.003 & C-12(p,x)C-11 & 5700.00--5700.00 & accepted & 1 & 0 \\
C0221.002 & C-12(p,x)C-11 & 3000.00--6000.00 & accepted & 5 & 0 \\
C0221.003 & C-12(p,x)C-11 & 4100.00--4100.00 & accepted & 1 & 0 \\
C0231.002 & C-12(p,x)C-11 & 28000.00--28000.00 & accepted & 1 & 0 \\
C1911.002 & C-12(p,x)C-11 & 800.00--800.00 & accepted & 1 & 0 \\
C1911.003 & C-12(p,x)C-11 & 800.00--800.00 & accepted & 2 & 0 \\
C2341.002 & C-12(p,x)C-11 & 300000.00--300000.00 & accepted & 1 & 0 \\
C2359.002 & C-12(p,x)C-11 & 283.00--387.00 & accepted & 2 & 0 \\
C2359.003 & C-12(p,x)C-11 & 400.00--400.00 & accepted & 1 & 0 \\
D0128.005 & C-12(p,x)C-11 & 150.00--150.00 & accepted & 1 & 1 \\
E2449.005 & C-12(p,x)C-11 & 32.50--69.60 & accepted & 4 & 4 \\
E2501.002 & C-12(p,x)C-11 & 10.53--67.61 & accepted & 35 & 32 \\
F1218.002 & C-12(p,x)C-11 & 3650.00--3650.00 & accepted & 1 & 0 \\
O0415.002 & C-12(p,x)C-11 & 7600.00--7600.00 & accepted & 1 & 0 \\
O0429.002 & C-12(p,x)C-11 & 1000.00--1000.00 & accepted & 1 & 0 \\
O2553.003 & C-12(p,x)C-11 & 40.70--215.10 & accepted & 8 & 5 \\
O2582.002 & C-12(p,x)C-11 & 23.60--189.30 & accepted & 43 & 35 \\
P0051.002 & C-12(p,x)C-11 & 27.60--41.80 & accepted & 31 & 1 \\
P0051.003 & C-12(p,x)C-11 & 18.00--31.80 & accepted & 36 & 0 \\
C0062.012 & N-14(p,x)N-13 & 155.00--5700.00 & accepted & 2 & 0 \\
C0062.017 & N-14(p,x)N-13 & 18.50--156.00 & accepted & 16 & 13 \\
C0202.004 & N-14(p,x)N-13 & 12.62--29.80 & accepted & 35 & 11 \\
C0220.005 & N-14(p,x)N-13 & 5700.00--5700.00 & accepted & 1 & 0 \\
C0244.002 & N-14(p,x)N-13 & 400.00--3000.00 & accepted & 4 & 0 \\
C1979.002.1 & N-14(p,x)N-13 & 400.00--6200.00 & accepted & 6 & 0 \\
C1979.002.2 & N-14(p,x)N-13 & 400.00--6200.00 & accepted & 5 & 0 \\
C1979.002.3 & N-14(p,x)N-13 & 720.00--720.00 & accepted & 1 & 0 \\
D4128.003 & N-14(p,x)N-13 & 11.40--19.20 & accepted & 10 & 10 \\
D4145.002 & N-14(p,x)N-13 & 19.15--25.62 & accepted & 4 & 4 \\
O2067.027 & N-14(p,x)N-13 & 145.00--145.00 & accepted & 1 & 1 \\
O2134.003 & N-14(p,x)N-13 & 50.00--50.00 & accepted & 1 & 1 \\
O2582.007 & N-14(p,x)N-13 & 22.70--189.10 & accepted & 43 & 36 \\
C0061.019 & O-16(p,x)C-11 & 155.00--155.00 & accepted & 1 & 0 \\
C0208.003 & O-16(p,x)C-11 & 362.00--362.00 & accepted & 1 & 0 \\
C0220.011 & O-16(p,x)C-11 & 5700.00--5700.00 & accepted & 1 & 0 \\
C0223.014 & O-16(p,x)C-11 & 146.00--146.00 & accepted & 1 & 1 \\
D1061.002 & O-16(p,x)C-11 & 22.90--56.90 & accepted & 16 & 15 \\
E2449.003 & O-16(p,x)C-11 & 31.60--69.10 & accepted & 7 & 7 \\
E2568.004 & O-16(p,x)C-11 & 29.30--69.93 & pending verification & 133 & 133 \\
O0065.004 & O-16(p,x)C-11 & 18.60--27.80 & accepted & 4 & 0 \\
O2134.006 & O-16(p,x)C-11 & 50.00--50.00 & accepted & 1 & 1 \\
O2134.007 & O-16(p,x)C-11 & 100.00--153.00 & accepted & 2 & 1 \\
O2576.003 & O-16(p,x)C-11 & 5.00--15.80 & accepted & 13 & 0 \\
O2582.003 & O-16(p,x)C-11 & 20.20--189.10 & accepted & 42 & 33 \\
B0157.002 & O-16(p,x)N-13 & 6.70--9.27 & accepted & 29 & 29 \\
C0061.018 & O-16(p,x)N-13 & 155.00--155.00 & accepted & 1 & 0 \\
C0202.002 & O-16(p,x)N-13 & 6.66--15.77 & accepted & 37 & 37 \\
C0208.005.2 & O-16(p,x)N-13 & 362.00--362.00 & accepted & 1 & 0 \\
C0220.010 & O-16(p,x)N-13 & 5700.00--5700.00 & accepted & 1 & 0 \\
C0223.013 & O-16(p,x)N-13 & 146.00--146.00 & accepted & 1 & 1 \\
C0998.002 & O-16(p,x)N-13 & 12.23--17.79 & accepted & 355 & 355 \\
D0138.003 & O-16(p,x)N-13 & 13.10--13.10 & accepted & 1 & 1 \\
D1061.003 & O-16(p,x)N-13 & 12.00--56.90 & accepted & 17 & 17 \\
E2449.004 & O-16(p,x)N-13 & 7.20--69.10 & accepted & 9 & 8 \\
E2568.003 & O-16(p,x)N-13 & 5.65--69.93 & pending verification & 167 & 167 \\
O0065.005 & O-16(p,x)N-13 & 15.60--27.80 & accepted & 5 & 0 \\
O2134.008 & O-16(p,x)N-13 & 153.30--153.30 & accepted & 1 & 0 \\
O2576.004 & O-16(p,x)N-13 & 6.50--16.40 & accepted & 13 & 13 \\
O2582.004 & O-16(p,x)N-13 & 20.20--189.10 & accepted & 42 & 35 \\
P0045.003 & O-16(p,x)N-13 & 9.73--15.00 & accepted & 32 & 4 \\
P0051.004 & O-16(p,x)N-13 & 6.22--15.60 & accepted & 60 & 1 \\
R0043.002 & O-16(p,x)N-13 & 6.51--7.71 & accepted & 11 & 11 \\
A0316.008 & O-16(p,x)O-15 & 16.70--32.82 & accepted & 15 & 15 \\
C0062.013 & O-16(p,x)O-15 & 155.00--5700.00 & accepted & 2 & 0 \\
C0062.016 & O-16(p,x)O-15 & 19.50--156.00 & accepted & 18 & 15 \\
C0208.004.2 & O-16(p,x)O-15 & 362.00--362.00 & accepted & 1 & 0 \\
C0220.008 & O-16(p,x)O-15 & 5700.00--5700.00 & accepted & 1 & 0 \\
C0223.012 & O-16(p,x)O-15 & 146.00--146.00 & accepted & 1 & 1 \\
D1061.004 & O-16(p,x)O-15 & 12.00--56.90 & accepted & 17 & 16 \\
E2449.002 & O-16(p,x)O-15 & 7.20--69.10 & accepted & 9 & 8 \\
E2568.002 & O-16(p,x)O-15 & 11.98--69.93 & pending verification & 161 & 158 \\
O0065.006 & O-16(p,x)O-15 & 18.60--27.80 & accepted & 4 & 0 \\
O2553.006 & O-16(p,x)O-15 & 39.40--217.00 & accepted & 5 & 4 \\
O2582.005 & O-16(p,x)O-15 & 20.20--189.10 & accepted & 42 & 35 \\
\bottomrule
\end{longtable}

\endgroup

\section{Point-level curation audit}

Table~\ref{tab:point-curation} partitions all selected EXFOR points by their
first applicable selection decision.  ``No error'' denotes the absence of a
positive reported cross-section uncertainty; ``Range'' denotes energies
outside \SIrange{5}{150}{MeV}; and ``Pending'' denotes points belonging to a
dataset pending verification.

\begin{table}[ht]
\centering\small
\caption{Point-level curation outcome by channel.}
\label{tab:point-curation}
\resizebox{\textwidth}{!}{\begin{tabular}{lrrrrrr}
\toprule
Channel & Used & No error & Range & Threshold & Nonpositive & Pending \\
\midrule
C-12(p,x)C-11 & 94 & 104 & 55 & 7 & 0 & 59 \\
O-16(p,x)O-15 & 94 & 6 & 13 & 2 & 0 & 161 \\
O-16(p,x)C-11 & 58 & 1 & 11 & 19 & 0 & 133 \\
N-14(p,x)N-13 & 76 & 26 & 27 & 0 & 0 & 0 \\
O-16(p,x)N-13 & 512 & 92 & 11 & 0 & 1 & 167 \\
\bottomrule
\end{tabular}
}
\end{table}

\section{Numerical configuration and reproducibility}

The authoritative curation and fit settings are
\path{config/xsection_curation.toml} and
\path{config/xsection_fit.toml}.  The fit metadata record their SHA-256
digests, the random seed, the spline knots, the candidate curvature penalties,
the candidate campaign spreads, the replica count, and the output grids.
Dense curves, common-grid tables, all replicas, representative curves, and
distance histograms are stored under \path{data/xsections/fits}.

The analysis and document are regenerated with
\begin{verbatim}
python3 tools/build_xsection_fit.py
\end{verbatim}
The build regenerates the normalized input, both curation audits, fitted
ensembles, tables, figures, validation checks, and this document.

\end{document}
